\begin{frame}[t, fragile]{Types and type classes in Haskell.}

In Haskell, there is a distinction between types -- or data types if you prefer, and type classes. \\

\textbullet{} Types \\

Haskell has a number of built-in or primitive types. These include types
such as \mintinline[fontsize=\normalsize]{haskell}{Char},
\mintinline[fontsize=\normalsize]{haskell}{String},
\mintinline[fontsize=\normalsize]{haskell}{Int}, \mintinline[fontsize=\normalsize]{haskell}{Double},
\mintinline[fontsize=\normalsize]{haskell}{Float} and the like. \\

\textbullet{} Type classes \\

In Haskell, a type class isn't actually a type. Rather, a type class is essentially a collection of operations which any new type that derives from it can implement. So a type class is not a type which can be used directly, but instead, it acts at the basis for creating new types which \textit{can} be used directly. All of the new types which are instances of a particular type class, will support the same particular set of operations which are defined in the type class.

Put another way, a type class defines a common interface that many different types can implement.

\end{frame}